% !TEX root = ../main.tex
\chapter{\texttt{example} と \texttt{theorem}}
\label{chap:ch06_example_theorem}
\BookIndex{theorem}{\texttt{theorem}}
\BookIndex{example}{\texttt{example}}

\begin{chapterabstract}
本章では、証明に名前を付けるかどうかの違い、すなわち \texttt{example} と \texttt{theorem} を扱います。
前章で定義した \texttt{DiscountSpec} を使い、\texttt{theorem member\_hasDiscount : DiscountSpec true := by rfl} に到達します。
ここまで進むと、定理を見ただけで、仕様名の展開から \texttt{rfl} で閉じるまでの流れを追えるようになります。
\end{chapterabstract}

\begin{learninggoals}
\item \texttt{example} と \texttt{theorem} の違いを、名前の有無と再利用の可否によって説明できます。
\item 仕様名に対する定理を書き、\texttt{rfl} で証明できます。
\item 定理のゴールが、仕様名の展開を通じてどのように閉じるかを表で追えます。
\end{learninggoals}

\section{問題設定}

前章までで、仕様に名前を付け、その意味を展開できるようになりました。
本章では、その仕様が成り立つことを証明として保存します。
証明の書き方には、名前を付けない \texttt{example} と、名前を付ける \texttt{theorem} があります。
後からその証明を参照するかどうかが、使い分けの基準です。

図\ref{fig:ch06-spec-proof-flow}は、関数の振る舞いを命題として切り出し、その命題に証明を与えるまでを三段階に分けて示します。

\FigurePlaceholder{fig-ch06-spec-proof-flow.png}{0.92}{関数・命題・証明}{fig:ch06-spec-proof-flow}

\section{名前付きの定理を書く}

仕様名に対する証明に名前を付けると、後から別の証明や設計文書で参照できます。
\texttt{DiscountSpec true} が成り立つことを、名前付きの定理として書きます。

\begin{listing}[htbp]
\begin{minted}{lean4}
def hasDiscount (member : Bool) : Bool :=
  member

def DiscountSpec (member : Bool) : Prop :=
  hasDiscount member = true

theorem member_hasDiscount : DiscountSpec true := by
  rfl

-- これは証明できない。
-- DiscountSpec false は展開すると false = true になる。
-- theorem nonmember_hasDiscount : DiscountSpec false := by
--   rfl
\end{minted}
\caption{\texttt{member\_hasDiscount}}
\label{lst:ch06_theorem}
\end{listing}

コード\ref{lst:ch06_theorem}の \texttt{member\_hasDiscount} は、\texttt{DiscountSpec true} が成り立つことを示す定理です。
\texttt{rfl} で閉じる理由は、ゴールの展開を順に追うと分かります。
一方、コメントアウトした \texttt{nonmember\_hasDiscount} は、\texttt{DiscountSpec false} が \texttt{false = true} に展開されるため証明できません。

\begin{table}[htbp]
\centering
\caption{\texttt{member\_hasDiscount} の展開}
\label{tab:ch06-goal}
\begin{tabular}{ll}
\toprule
読む段階 & 内容 \\
\midrule
定理名 & \texttt{member\_hasDiscount} \\
証明したいこと & \texttt{DiscountSpec true} \\
\texttt{DiscountSpec} を展開 & \texttt{hasDiscount true = true} \\
\texttt{hasDiscount} を展開 & \texttt{true = true} \\
使う証明 & \texttt{rfl} \\
閉じる理由 & 左右が同じ \\
\bottomrule
\end{tabular}
\end{table}

\section{名前のない確認は \texttt{example} で書く}

その場限りの確認には、名前を付けない \texttt{example} が適しています。
ここでは、税込金額の具体例を \texttt{example} で確かめます。

\begin{listing}[htbp]
\begin{minted}{lean4}
def addTax (price tax : Nat) : Nat :=
  price + tax

def TaxSpec (price tax total : Nat) : Prop :=
  total = addTax price tax

example : TaxSpec 1000 100 1100 := by
  rfl
\end{minted}
\caption{\texttt{TaxSpec} の具体例}
\label{lst:ch06_example}
\end{listing}

コード\ref{lst:ch06_example}の \texttt{TaxSpec 1000 100 1100} は、\texttt{1100 = addTax 1000 100}、すなわち \texttt{1100 = 1000 + 100} に展開され、\texttt{1100 = 1100} になります。
この \texttt{example} には名前がないため、後から別の証明で参照できません。

\section{本章のまとめ}

\begin{summarybox}
要点と記法
\begin{itemize}
\item \texttt{theorem 名前 : P := by ...}：名前付きの証明
\item \texttt{example : P := by ...}：名前のない証明
\end{itemize}
\end{summarybox}
